\documentclass{article}
\usepackage[T2A]{fontenc}
\usepackage[utf8]{inputenc}   
\usepackage[english, russian]{babel}

% Set page size and margins
% Replace `letterpaper' with`a4paper' for UK/EU standard size
\usepackage[a4paper,top=2cm,bottom=2cm,left=2cm,right=2cm,marginparwidth=1.75cm]{geometry}

\usepackage{amsmath}
\usepackage{graphicx}
\usepackage[colorlinks=true, allcolors=blue]{hyperref}
\usepackage{amsfonts}
\usepackage{amssymb}
% \usepackage[left=1cm,right=1cm,top=1cm,bottom=1cm]{geometry}
\usepackage{hyperref}
\usepackage{seqsplit}
\usepackage[dvipsnames]{xcolor}
\usepackage{enumitem}
\usepackage{algorithm}
\usepackage{algpseudocode}
\usepackage{algorithmicx}
\usepackage{mathalfa}
\usepackage{mathrsfs}
\usepackage{dsfont}
\usepackage{caption,subcaption}
\usepackage{wrapfig}
\usepackage[stable]{footmisc}
\usepackage{indentfirst}
\usepackage{rotating}
\usepackage{pdflscape}
\usepackage{listings}
\usepackage{amsmath}
\usepackage{graphicx}
\usepackage[colorlinks=true, allcolors=blue]{hyperref}
\usepackage{amsfonts}
\usepackage{amssymb}
% \usepackage[left=1cm,right=1cm,top=1cm,bottom=1cm]{geometry}
\usepackage{hyperref}
\usepackage{seqsplit}
\usepackage[dvipsnames]{xcolor}
\usepackage{enumitem}
\usepackage{algorithm}
\usepackage{algpseudocode}
\usepackage{algorithmicx}
\usepackage{mathalfa}
\usepackage{mathrsfs}
\usepackage{dsfont}
\usepackage{caption,subcaption}
\usepackage{wrapfig}
\usepackage[stable]{footmisc}
\usepackage{indentfirst}
\usepackage{rotating}
\usepackage{pdflscape}
\usepackage{graphicx}

\usepackage{minted}

\usepackage{MnSymbol,wasysym}
\usepackage{tikz}

\usepackage{minted}

\usepackage{MnSymbol,wasysym}

\title{Домашнее задание 19}
\author{Лиза Фоменко}
\date{}

\begin{document}

\maketitle


\section*{\textbf{Задание 1}} \textbf{\textit{Найдите сложность построения кучи последним методом, который был (частично) обсужден на лекции. Для простоты будем считать, что в куче $n = 2^k - 1$ элементов.
На первом шаге вызывается процедура {heapify} для подкуч размера 3, на втором размера 7 и так далее. Считаем, что {heapify} работает за $O(h)$, где $h$ - высота кучи, для которой он вызван.}}

\textit{heapify} работает за $O(h)$, где $h$ - высота кучи. 

Алгоритм построения кучи:

1) строим рандомное дерево из наших $n = 2^k - 1$ элементов.

2) вызываем операцию \textit{heapify} для всех поддеревьев, начиная с наименьших, постепенно увеличиваем размер подкучек.

Сначала пукнт 1: построение дерева из $n = 2^k - 1$ элементов. Построение такого дерева займет $O(n)$.

Далее пункт 2: начнем с лекции

$$\sum_{i = 1}^{\lg_2 n}  2^{\lg_2 n-i}* heapify(i) = \sum_{i = 1}^{\lg_2 n}  \frac{n}{2^i}*O(i) = Cn   \sum_{i = 1}^{\lg_2 n}  \frac{i}{2^i}$$

Для взятия этого ряда вспомним, что $\frac{d}{dx} x^{n} = nx^{n-1}$

$\sum_{i=1}^{\infty} x^{n} = \frac{1}{1-x}$, где $|x| < 1$

тогда

$\sum_{i=1}^{\infty} nx^{n-1} = \frac{1}{(1-x)^2}$

и

$\sum_{i=1}^{\infty} nx^{n} = \frac{x}{(1-x)^2}$

У нас то же самое:

$$\sum_{i = 1}^{\infty}  \frac{i}{2^i} = \frac{\frac{1}{2}}{(1-\frac{1}{2})^2} = 2$$

То есть ряд $\sum_{i = 1}^{\infty}  \frac{i}{2^i}$ сходится к 2. Наш ряд немного другой, $\sum_{i = 1}^{\lg_2 n}  \frac{i}{2^i}$, но он также будет сходиться к константе по признакам сходимости (в нем просто меньше слагаемых). 

Тогда получаем, что 

$$Cn   \sum_{i = 1}^{\lg_2 n}  \frac{i}{2^i} \leq 2Cn = O(n)$$

то есть второй пункт их алгоритма работает за $O(n)$

Тогда весь алгоритм, вместе с построением рандомного дерева, работает за $O(n)$.


\bigskip

\section*{\textbf{Задание 2}} \textbf{\textit{Приведите алгоритм добавления элемента в уже существующую кучу на максимум из $n$ элементов. Докажите его корректность и оцените асимптотику.}}

Как добавить элемент в кучу?

\begin{verbatim}
heap  # уже есть куча из n элементов

def add_to_heap(element):
    heap[n+1] = element  # если куча дана массивом, ставим эелемент на ласт место
    # если куча реализована как дерево, 
    # наш элемент займет самое первое (левое) свободное место в последнем ряду узлов
    
\end{verbatim}

Смысл: сначала поставим элемент на первое свободное место в куче: в массиве из $n$ элементов это $n+1$ позиция, если куча дана в виде дерева бинарного, то если последний уровень не заполнен, ставим элемент на последний уровень самым левым (если уровень заполнен, то ставим его первым (самым левым) на новый уровень)

За какое время мы добавим? если массив, то $O(1)$. Если мы храним кучу в виде дерева, то мы можем вычислить его место в дереве за $O(lgN)$. 

Что мы теперь значем о куче? Любой родитель больше своих потомков (так как добавляли элемент в кучу с выполненным свойством), кроме, быть может, родителя только что добавленного элемента (мы за этим не следили). Тогда нам необходимо исправить только восходящую ветвь потомков вновь добавленного узла, из $lgN$. Мы будем делать операцию $change$, которая работает за $O(1)$.

\begin{verbatim}
def change(node):
    if node.value > node.parent.value:
        node.value, node.parent.value = node.parent.value, node.value 
\end{verbatim}

Так как у нас нет ограничений на добавляемый элемент, он в теории может быть больше любого своего родителя; тогда нам потребуется сделать операцию $change$ столько раз, сколько предков у нашего листа, а их $lgN$. Тогда восстановление свойства кучи займет $O(lgN)$

Тогда обе операции занимают $O(lgN)$ операций.
\bigskip


\section*{\textbf{Задание 3}} \textbf{\textit{На вход задачи поступают $n$ точек на плоскости. Найдите точку, для которой сумма евклидовых расстояний от нее до входных минимальна.
}}

Точка, минимизирующая сумму евклидовых расстояний, это геометрическая медиана (точка Ферма)

$\sum_{i=1}^n \sqrt{(x_1 - x_{i_1})^2 + (x_2 - x_{i_2})^2}$ -> $min$

найдем производные по $x_1$ и $x_2$

$$\frac{d}{dx_1} \sum_{i=1}^n \sqrt{(x_1 - x_{i_1})^2 + (x_2 - x_{i_2})^2} = \sum_{i=1}^n \frac{d}{dx_1} \sqrt{(x_1 - x_{i_1})^2 + (x_2 - x_{i_2})^2} = \sum_{i=1}^n \frac{2(x_1 - x_{i_1})}{2 \sqrt{(x_1 - x_{i_1})^2 + (x_2 - x_{i_2})^2}} = $$ $$ = \sum_{i=1}^n \frac{x_1 - x_{i_1}}{\sqrt{(x_1 - x_{i_1})^2 + (x_2 - x_{i_2})^2}} = 0$$

$$\frac{d}{dx_2} \sum_{i=1}^n \sqrt{(x_1 - x_{i_1})^2 + (x_2 - x_{i_2})^2} = \sum_{i=1}^n \frac{d}{dx_2} \sqrt{(x_1 - x_{i_1})^2 + (x_2 - x_{i_2})^2} = \sum_{i=1}^n \frac{2(x_2 - x_{2_1})}{2 \sqrt{(x_1 - x_{i_1})^2 + (x_2 - x_{i_2})^2}} = $$ $$ = \sum_{i=1}^n \frac{x_2 - x_{2_1}}{\sqrt{(x_1 - x_{i_1})^2 + (x_2 - x_{i_2})^2}} = 0$$

То есть получаем систему:

\begin{equation}
\left\{
\begin{array}{l}
\sum_{i=1}^n \frac{x_1 - x_{i_1}}{\sqrt{(x_1 - x_{i_1})^2 + (x_2 - x_{i_2})^2}} = 0  \\
\sum_{i=1}^n \frac{x_2 - x_{2_1}}{\sqrt{(x_1 - x_{i_1})^2 + (x_2 - x_{i_2})^2}} = 0 \quad
\end{array}
\right.
\end{equation}

Вообще, аналитически она не решается. Кажется для ее нахождения используют приближенные методы

\textbf{ОКЕЙ НУЖНО ДЛЯ КВАДРАТОВ РАССТОЯНИЙ, ПЕРЕРЕШИВАЮ НИЖЕ}

Тут уже ищем центр масс)

$\sum_{i=1}^n (x_1 - x_{i_1})^2 + (x_2 - x_{i_2})^2$ -> $min$

найдем производные по $x_1$ и $x_2$

$$\frac{d}{dx_1} \sum_{i=1}^n (x_1 - x_{i_1})^2 + (x_2 - x_{i_2})^2 = \sum_{i=1}^n \frac{d}{dx_1} ((x_1 - x_{i_1})^2 + (x_2 - x_{i_2})^2) = \sum_{i=1}^n 2(x_1 - x_{i_1}) = 2\sum_{i=1}^n (x_1 - x_{i_1}) = 2 (nx_1 - \sum_{i=1}^n x_{i_1}) = 0$$

откуда $$x_1 = \frac{\sum_{i=1}^n x_{i_1}}{n}$$

$$\frac{d}{dx_2} \sum_{i=1}^n (x_1 - x_{i_1})^2 + (x_2 - x_{i_2})^2 = \sum_{i=1}^n \frac{d}{dx_2} ((x_1 - x_{i_1})^2 + (x_2 - x_{i_2})^2) = \sum_{i=1}^n 2(x_2 - x_{i_2}) = 2\sum_{i=1}^n (x_2 - x_{i_2}) = 2 (nx_2 - \sum_{i=1}^n x_{i_2}) = 0$$

откуда $$x_2 = \frac{\sum_{i=1}^n x_{i_2}}{n}$$

Тогда ответ: искомая точка  $$ (x_1, x_2) = \left( \frac{\sum_{i=1}^n x_{i_1}}{n}, \frac{\sum_{i=1}^n x_{i_2}}{n} \right)$$ 
\bigskip

\section*{\textbf{Задание 4}} \textbf{\textit{Примените МНК к задаче нахождения параболы по набору из $n$ точек на плоскости. Достаточно построить систему линейных уравнений, решение которой даст ответ задачи.}}

$\{x_i, y_i\}_{i=1}^{n}$

$f(x) = ax^2 + bx + c$

По МНК:

$E = \sum_{i=1}^n e_i = \sum_{i=1}^n (y_i - f(x))= \sum_{i=1}^n (y_i - ax_i^2 - bx_i- c)^2$ -> $min$

Найдем производные по неизвестным параметрам $a, b, c$

$\frac{d}{da} \sum_{i=1}^n (y_i - ax_i^2 - bx_i- c)^2 =  \sum_{i=1}^n \frac{d}{da}(y_i - ax_i^2 - bx_i- c)^2 = \sum_{i=1}^n -2x_i^2(y_i - ax_i^2 - bx_i- c) = 0$

$\frac{d}{db} \sum_{i=1}^n (y_i - ax_i^2 - bx_i- c)^2 =  \sum_{i=1}^n \frac{d}{db}(y_i - ax_i^2 - bx_i- c)^2 = \sum_{i=1}^n -2x_i(y_i - ax_i^2 - bx_i- c) = 0$

$\frac{d}{dc} \sum_{i=1}^n (y_i - ax_i^2 - bx_i- c)^2 =  \sum_{i=1}^n \frac{d}{dc}(y_i - ax_i^2 - bx_i- c)^2 = \sum_{i=1}^n -2(y_i - ax_i^2 - bx_i- c) = 0$

\begin{equation}
\left\{
\begin{array}{l}
\sum_{i=1}^n x_i^2(y_i - ax_i^2 - bx_i- c) = 0  \\
\sum_{i=1}^n x_i(y_i - ax_i^2 - bx_i- c) = 0 \\
\sum_{i=1}^n (y_i - ax_i^2 - bx_i- c) = 0
\end{array}
\right.
\end{equation}

можно еще довести до вида:

\begin{equation}
\left\{
\begin{array}{l}
 \sum_{i=1}^n x_i^2(ax_i^2 + bx_i + c) = \sum_{i=1}^n x_i^2y_i  \\
\sum_{i=1}^n x_i(ax_i^2 + bx_i + c) = \sum_{i=1}^n x_iy_i \\
 \sum_{i=1}^n (ax_i^2 + bx_i + c) = \sum_{i=1}^n y_i
\end{array}
\right.
\end{equation}




\bigskip

\section*{\textbf{Задание 5}} \textbf{\textit{Сигнал имеет вид $x(t) = a \sin(t)$. На вход алгоритма поступает набор данных $\{ (t_i, x_i) \}_{i=1}^n$. Примените МНК для нахождения значения параметра $\hat{a}$, при котором достигается минимум квадратичной функции ошибки.}}

$\{ (t_i, x_i) \}_{i=1}^n$

$f(t) = a \sin(t)$

$\sum_{i=1}^n (x_i - f(t))^2 = \sum_{i=1}^n (x_i - a \sin(t_i))^2 $ -> $min$

Ищем $\hat{a}$.

$\frac{d}{da} \sum_{i=1}^n (x_i - a \sin(t_i))^2 = \sum_{i=1}^n \frac{d}{da}(x_i - a \sin(t_i))^2 = \sum_{i=1}^n -2\sin(t_i)(x_i - a \sin(t_i)) = 0$

Тогда

$\sum_{i=1}^n \sin(t_i)(x_i - a \sin(t_i)) = 0$

$\sum_{i=1}^n x_i\sin(t_i) - a \sin(t_i)^2 = 0$

$\sum_{i=1}^n x_i\sin(t_i) = \sum_{i=1}^n a \sin(t_i)^2 = a\sum_{i=1}^n \sin(t_i)^2$

Тогда 

$$\hat{a} = \frac{\sum_{i=1}^n x_i\sin(t_i)}{\sum_{i=1}^n \sin(t_i)^2}$$

\bigskip

\section*{\textbf{Задание 6}} \textbf{\textit{На вход задачи поступает двумерный массив, состоящий из нулей и единиц и имеющий размеры $w \times h$. Единицами обозначены точки, лежащие на зашумленной линии. Найдите угловое и пространственное разрешение аккумуляторного массива, при которых асимптотическая сложность преобразования Хафа для прямых оценивается как $O(wh)$, действуя в предположении, что $h = \Theta(w)$ и $h = \Theta(n)$, где $n$ - число единиц в массиве.}}


\bigskip


\section*{\textbf{Задание 7}} \textbf{\textit{С помощью алгоритма $RANSAC$ производится детектирование сфер в трехмерном пространстве. Ему на вход поступают данные со стереокамеры в формате множества трехмерных точек $\{ (x_i, y_i, z_i) \}_{i=1}^n$. От алгоритма требуется работать $M$ раз в секунду. На обработку каждого облака точек тратится одинаковое время. Тестирование одной гипотезы в процессе выполнения алгоритма (нахождение подвыборки, построение модели и подсчет числа подхдящих под модель точек) занимает $t$ секунд ($t < 1$). Найдите вероятность того, что алгоритм будет давать корректный ответ в течение одной секунды подряд (то есть $M$ раз), пренебрегая всеми операциями, кроме тестирования гипотез. Доля инлаеров (не шумовых точек) составляет $\alpha$.}}

Так, нам нужно выбрать 4 (для сферы нужно 4 точки) точки-инлайера. Вероятность того, что мы выберем 4 инлайерных точки, примерно равна $\alpha ^4$. 

Вероятность, что гипотеза неудачна, $1 - \alpha^4$. 

Теперь такой расчет: у нас 1 секунда, в ней M запусков. Всего итераций в каждом запуске $\frac{1}{Mt}$, то есть за 1 секунду мы наблюдаем $\frac{1}{t}$ гипотез.

Тогда вероятность того, что хотя бы 1 гипотеза в запуске будет удачной, равна $1 - (1 - \alpha^4)^{\frac{1}{Mt}}$. 

Нам нужно M таких запусков подряд, это $$(1 - (1 - \alpha^4)^{\frac{1}{Mt}})^M$$






\bigskip

\end{document}